\chapter{Trabalhos Relacionados}\label{ch:related}

O aumento incessante do número de núcleos por processador em conjunto com a crescente heterogeneidade dos hardwares utilizados
acaba por expor limites e dores das arquiteturas dos sistemas operacionais monolíticos, onde um único kernel é responsável por
fazer o manejo de todas as estruturas de dados e controle da máquina.
Desde a década de 1990, uma linha de pesquisa propõe a alternativa de executar múltiplas instâncias de um mesmo kernel, ou
subconjunto de kernels, com diferentes configurações em uma mesma máquina física, onde todos funcionam de forma parecida com um sistema distribuído.

Essa maneira de tratar o SO hoje é conhecida como multikernel.
Este capítulo apresenta os trabalhos correlatos estudados, mostrando a evolução da arquitetura multikernel e conectando outros
trabalhos ao foco principal da pesquisa: desenvolver uma forma de orquestrar, monitorar e balancear cargas de trabalho
entre diferentes nós que utilizam o sistema multikernel de maneira escalável e previsível.

Este capítulo organiza os trabalhos correlatos em quatro seções.
A Seção~\ref{sec:rel-multikernel} apresenta a evolução das arquiteturas com múltiplos kernels, desde o início com os sistemas para
multiprocessadores de memória compartilhada até multikernels reduzidos utilizados em ambientes HPC.
A Seção~\ref{sec:rel-orquestracao} discute como esses sistemas e orquestradores decidem a alocação de recursos
e o balanceamento de carga dentro da própria arfquitetura e em ambientes distribuídos.
A Seção~\ref{sec:rel-observabilidade} aborda monitoramento, o qual é o critério para a tomada das decisões de orquestração.
Já a Seção~\ref{sec:rel-disponibilidade} trata sobre isolamento e disponibilidade dessas instâncias e quais seus objetivos em geral.
Por fim, a Seção~\ref{sec:sumario-e-trabalhos-revisados} sintetiza rapidamente tudo o que foi discutido.

% ---------------------------------------------------------------------
\section{Arquiteturas com múltiplos kernels}
\label{sec:rel-multikernel}

A ideia de estruturar o SO como um conjunto de kernels independentes surgiu há algumas décadas.
O Hive \citep{chapin1995hive}, organiza o SO como um sistema distribuído interno de kernels, os quais são denominados células,
cada uma delas é responsável por um conjunto fixo de nós e opera como um SO independente, enquanto
cooperam para oferecer às aplicações uma imagem única e transparente de sistema.
A comunicação entre células ocorre por meio de \textit{remote procedure calls} (RPC) e por leitura direta da memória de outras células.

Em vez de reescrever o SO por completo, Disco \citep{bugnion1997disco} ataca o problema inserindo um pequeno monitor de máquinas virtuais
que executa várias cópias de um SO na mesma máquina, comunicando-se por protocolos distribuídos comuns
como NFS e TCP/IP.
A abordagem permite executar sistemas operacionais especializados lado a lado com sistemas comerciais.

O Cellular Disco \citep{govil2000cellular} estende a proposta dividindo o próprio monitor em partes menores, a fim de conter falhas
de hardware, transformando o multiprocessador no que seria um cluster virtual.

Já o sistema Corey \citep{boydwickizer2008corey} parte do princípio de que as aplicações devem ser responsáveis por controlar o compartilhamento.
Por padrão, cada estrutura de dados do kernel é usada por um único núcleo, e abstrações como \textit{address ranges}, \textit{kernel cores}
e \textit{shares} permitem à aplicação indicar explicitamente o que deve ser compartilhado.
Por mais que execute um único kernel, o sistema Corey foi uma base importante para trabalhos futuros como Barrelfish.

Enquanto isso, o FOS \citep{wentzlaff2009fos} propõe fatorar os serviços do SO em frotas de servidores distribuídos
espacialmente pelos núcleos, permitindo a comunicação exclusivamente por mensagens, isso foi motivado pela constatação de que
o alocador de páginas físicas do Linux não escalava para além de oito núcleos.
Além disso, essa ideia da comunicação forte por mensagens obteve grande aceitação por futuras propostas de multikernels como
a implementação em \citet{wang2025multikernel} feita a partir do kernel do Linux.

A formulação moderna do conceito é apresentada por \citet{baumann2009multikernel}, que propõem tratar a máquina como uma
rede de núcleos independentes.
Responsáveis por desenvolver o sistema Barrelfish, onde é executado um kernel por núcleo e três princípios básicos são seguidos: comunicação entre núcleos sempre é
explícita, sua estrutura é neutra em relação ao hardware e o estado é replicado em vez de ser compartilhado, sendo mantido consistente por processos monitores.
Em um trabalho anterior do mesmo grupo, \citet{schupbach2008diversity} argumentam que a diversidade de hardware, nas formas de não
uniformidade de memória, heterogeneidade de núcleos e diversidade entre sistemas, será um desafio tão grande quanto a escalabilidade no futuro.

No ambiente HPC, a motivação principal é combinar o baixo overhead de um kernel reduzido, comumente denotado como kernel leve, com a compatibilidade do Linux.
FusedOS \citep{park2012fusedos} particiona os núcleos entre o Linux e uma biblioteca derivada de kernel leve CNK, executada
como um processo de usuário, responsável por gerenciar núcleos dedicados às aplicações.

O trabalho Interface for Heterogeneous Kernels (IHK) \citep{shimosawa2014ihk} generaliza essa estrutura em um modelo mestre-escravo,
no qual o Linux, como kernel mestre, reserva núcleos e memória, inicializa, consulta o estado, desliga e destrói instâncias
de kernels leves, que se comunicam por uma camada própria.

Ao comparar seis desses projetos, \citet{gerofi2016survey} optam por distinguir esses sistemas com particionamento estático de recursos,
de sistemas que permitem reparticionamento dinâmico, dando grande importância a este ponto.

Além disso, duas propostas exploram áreas complementares que merecem ser citadas.
PIKA \citep{beckmann2014pika} projeta uma pilha de rede específica para multikernels, dividindo em servidores que se comunicam por mensagens.
Enquanto MultiK \citep{kuo2019multik}, que apesar do nome não executa múltiplos kernels independentes, mas sim oferece a
cada aplicação uma cópia do código do kernel que contém apenas o necessário para ela executar.
Isso ocorre automaticamente durante a troca de contexto dos processos, tendo o objetivo de reduzir a superfície de ataque de
aplicações que não precisam de todos os elementos do sistema e dando ênfase na ideia de especializar o kernel por aplicação.

% ---------------------------------------------------------------------
\section{Orquestração e balanceamento de carga}
\label{sec:rel-orquestracao}

Decidir onde executar cada carga e como redistribuí-la conforme a demanda é um problema estudado há muito tempo em sistemas distribuídos.
\citet{casavant1988taxonomy} propõem uma taxonomia hierárquica de escalonamento em sistemas distribuídos que distingue abordagens estáticas de dinâmicas,
Além disso, as abordagens dinâmicas são classificadas quanto à distribuição da responsabilidade pelas decisões e à cooperação entre os nós.

\citet{shivaratri1992load} decompõem os algoritmos de distribuição de carga em quatro principais políticas, de transferência, decidir
quando mover uma tarefa, de seleção, qual tarefa mover, localização, para onde mover e de informação, quando coletar o estado dos nós.
Os algoritmos ainda são classificados naqueles que são iniciados pelo emissor, pelo receptor ou simétricos.

Dois resultados fundamentam políticas simples, mas eficientes e escaláveis de balanceamento.
Primeiro, \citet{mitzenmacher2001power} mostra que escolhendo a opção menos carregada entre apenas dois destinos aleatórios
reduz exponencialmente a carga máxima esperada em relação a simplesmente escolher aleatoriamente.
Em outro estudo, \citet{blumofe1999work} provam que o roubo de tarefas de maneira aleatorizada de um processador ativo para outro ocioso, executa computações
com tempo próximo do ótimo.

Ao revisar orquestradores como Kubernetes, Docker Swarm e Apache Mesos, \citet{swetha2025resource} identificam a escassez
de propostas de orquestração orientada pela utilização de recursos e comparam a adequação de cada ferramenta.
Na mesma linha, \citet{lu2025optimized} propõem uma implantação de contêineres em dois estágios: um algoritmo guloso distribui contêineres
entre máquinas virtuais com granularidade grossa enquanto outro algoritmo genético reparte os recursos dentro de cada uma das máquinas virtuais.

Agora em sistemas com múltiplos kernels, desde o início, era necessário saber como distribuir os recursos da máquina entre as instâncias.
No Hive, os mecanismos de compartilhamento de recursos são implementados pela cooperação entre as células, mas a política
é atribuída a um processo de usuário, o único com visão global do sistema, responsável por decidir de
quais células alocar memória e como escalonar processos \citep{chapin1995hive}.

No Disco, o monitor aplica políticas globais sobre as máquinas virtuais, movendo memória entre elas para evitar paginação
desnecessária e escalonando dinamicamente as vCPUs para equilibrar a carga \citep{bugnion1997disco}.
O Cellular Disco \citep{govil2000cellular} apresenta uma solução mais elaborada para sistemas multiprocessadores,
nele, cada processador mantém sua própria fila de processadores virtuais, combinando um
balanceador ocioso, que usa apenas informação local, com um balanceador periódico, que consulta uma árvore de carga global.

No Barrelfish, a base de conhecimento do sistema reúne informação pelo hardware e medições em tempo de execução, respondendo
a consultas de otimização, como a alocação de drivers a núcleos e a construção de árvores de mensagens adaptadas à máquina
\citep{baumann2009multikernel,schupbach2008diversity}.

Nos multikernels com foco em HPC, a gerência das instâncias tende a ser exposta como uma interface.
A IHK oferece operações para reservar núcleos e memória, inicializar, consultar, desligar e destruir instâncias
de kernels leves, porém, não define políticas sobre quando as utilizar \citep{shimosawa2014ihk}.

No projeto Hobbes, \citet{kocoloski2015composition} propõem a execução em enclaves distintos de um mesmo nó, com um
gerenciador de nó responsável pelo ciclo de vida dos enclaves e mecanismos de memória compartilhada entre os mesmos.

Mais recentemente, a proposta de suporte a múltiplas instâncias do Linux em uma mesma máquina \citep{wang2025multikernel}
utiliza a ideia de carregar as instâncias via \textit{kexec}, recebendo núcleos e memória que podem ser
realocados em tempo de execução por meio de uma \textit{device tree} e do mecanismo de \textit{hotplug}
\citep{wang2025lpc}.
A ferramenta kerf \citep{multikernel2025kerf} ficou responsável por gerenciar a criação, configuração, monitoramento e a
destruição dessas instâncias.

\section{Observabilidade e monitoramento}
\label{sec:rel-observabilidade}

Esta seção apresenta fundamentos de como observar, coletar dados e analisá-los em sistemas distribuídos, em seguida como
trabalhos aplicam técnicas parecidas em ambientes com múltiplos kernels.

\citet{mansourisamani1993monitoring} definem um modelo funcional de monitoramento de sistemas distribuídos dividido em
quatro atividades (geração, processamento, disseminação e apresentação de informação), tratadas como blocos que se combinam
em diferentes proporções e momentos da análise, sendo uma base essencial para os trabalhos que seguem.
Além disso, apontam a desatualização da informação como o principal problema existente dentro do monitoramento.

Em uma revisão recente voltada a arquitetura de microsserviços, \citet{faseeha2025observability} analisam ferramentas de
observabilidade em diferentes níveis, como de sistema, de serviço e de rede, organizadas em torno de três fontes que hoje já são comumente disseminadas:
\textit{metrics}, \textit{logs} e \textit{traces}.

Para métricas em clusters, Ganglia \citep{massie2004ganglia} adotam uma organização hierárquica: dentro de cada cluster,
os nós comunicam e trocam seus estados via \textit{multicast}.
Clusters distintos são federados por uma árvore de conexões \textit{peer-to-peer}, com armazenamento em bases de séries temporais.
O projeto privilegia baixa sobrecarga por nó e escalabilidade, ao custo de métricas agregadas e de uma menor granularidade temporal.

Agora, em um ambiente multikernel, a IHK concentra a gerência no kernel \textit{master} \citep{shimosawa2014ihk}, responsável por
tomar o controle do mapeamento da memória física e orquestrar a alocação e reserva de recursos para os kernels \textit{slave} do ambiente.

O LDMS \citep{agelastos2014ldms} mostram que é possível monitorar continuamente um sistema inteiro, com a frequência adequada, mesmo em ambientes HPC.
LDMS utiliza processos coletores nos nós e processos agregadores, que reúnem os dados.
Os autores quantificam a sobrecarga do monitoramento sobre as aplicações como baixa.

O Omni-Kernel \citep{kvalnes2015omni} medem o consumo de recursos por completo e o atribui a um processo,
uma máquina virtual ou uma transação.
Nele, o kernel é fatorado em componentes que se comunicam por mensagens, e escalonadores existem no meio do caminho,
responsáveis por controlar a ordem de processamento e contabilizar o consumo.

Em uma proposta de hardware, \citet{west2009coremonitors} dedicam um núcleo específico ao monitoramento, alimentado por detectores de
eventos distribuídos pelo processador, mostrando como uma instância de kernel poderia contabilizar e exportar seu próprio consumo.

Os próprios sistemas com múltiplos kernels também possuem sistemas de monitoramento e balanceamento por si só.
A base de conhecimento do Barrelfish incorpora medições em tempo de execução, como a latência de comunicação entre cada par de
núcleos \citep{baumann2009multikernel}, o balanceador periódico do Cellular Disco
depende de uma árvore de carga atualizada constantemente \citep{govil2000cellular} e o PIKA mede o tempo de serviço de cada
aplicação para orientar as escolhas realizadas \citep{beckmann2014pika}.
Porém, a informação é coletada e consumida numa única máquina, não sendo explorada a ideia do compartilhamento entre nós em um cluster.

% ---------------------------------------------------------------------
\section{Disponibilidade e Isolamento}
\label{sec:rel-disponibilidade}

Uma grande motivação para dividir o sistema em múltiplos kernels é a disponibilidade das aplicações.
No Hive \citep{chapin1995hive}, uma falha de hardware ou de software afeta apenas os processos que usam os recursos da célula prejudicada.
O sistema emprega uma proteção em cada página contra escritas de outras células, descarta de maneira preventiva as páginas que uma
célula que falhou poderia ter corrompido e possui detecção de falhas por heurísticas, cuja confirmação exige consenso das células sobreviventes.

O Disco trata cada máquina virtual como unidade de contenção de falhas de software \citep{bugnion1997disco}, e o Cellular Disco estende
essa contenção a falhas de hardware.
Em centenas de experimentos com injeção de falhas, a recuperação foi bem-sucedida em todos os casos, levando menos de meio segundo \citep{govil2000cellular}.

Além de conter falhas, a divisão em múltiplos kernels também é usada para isolar o desempenho das aplicações.
O Pisces \citep{ouyang2015pisces} permite a coexistência, em um mesmo nó, de múltiplos ambientes de SO independentes e isolados,
denominados enclaves, cada um executando um co-kernel leve, capaz de gerenciar uma partição de recursos de hardware, tratando o
isolamento de desempenho como requisito principal e tendo seu foco em HPC

Com ênfase na parte de segurança, pesquisas como BlackBox \citep{vanthof2022blackbox} buscam proteger a confidencialidade e a
integridade de contêineres sem confiar no SO, por meio de um monitor que cria espaços de endereçamento físico protegidos
para cada contêiner.
Enquanto o MultiK \citep{kuo2019multik} garante maior segurança reduzindo a superfície de ataque ao remover mais de 90\% do
código do kernel visível a cada aplicação que está sendo executada.
Esses trabalhos, por mais que não executados sob a arquitetura multikernel, ressoam fortemente com sua filosofia.


\section{Sumário e trabalhos revisados}\label{sec:sumario-e-trabalhos-revisados}

A Tabela~\ref{tab:relacionados_ibn} apresenta os trabalhos selecionados e descritos, além de apresentar
um resumo sobre as principais contribuições e foco de cada um deles.

\begin{small} % Reduz o tamanho da fonte para que a tabela caiba melhor na página
    \begin{longtable}{p{3cm} p{4cm} p{4cm} p{3cm}}

% CABEÇALHO DA TABELA (Repetido em cada nova página)
        \caption{Resumo dos Trabalhos Relacionados.}
        \label{tab:relacionados_ibn} \\
        \toprule
        \textbf{Referência} & \textbf{Categoria} & \textbf{Contribuição principal} & \textbf{Relação com este trabalho} \\
        \midrule
        \endfirsthead % <--- Define que a parte acima é EXCLUSIVA da primeira página
% CABEÇALHO DA PRIMEIRA PÁGINA

% --- CABEÇALHO DAS PÁGINAS SEGUINTES (Sem o \label) ---
        \caption*{Resumo dos Trabalhos Relacionados. (Continuação)} \\
        \toprule
        \textbf{Referência} & \textbf{Categoria} & \textbf{Contribuição principal} & \textbf{Relação com este trabalho} \\
        \midrule
        \endhead % <--- Define o cabeçalho repetido nas páginas seguintes

% --- RODAPÉ INTERMEDIÁRIO (Páginas anteriores à última) ---
        \bottomrule
        \multicolumn{4}{r}{\footnotesize Continua na próxima página} \\
        \endfoot

% --- RODAPÉ FINAL (Apenas na última página) ---
        \bottomrule
        \endlastfoot

% --- Arquiteturas com múltiplos kernels ---
        \citet{chapin1995hive} & Multikernel & Hive: células com kernels independentes, contenção de falhas e política global no Wax & Origem da separação entre mecanismo e política \\
        \midrule
        \citet{bugnion1997disco} & Multikernel & Disco: múltiplas cópias de SO sobre monitor de máquinas virtuais & Políticas globais em instâncias independentes \\
        \midrule
        \citet{govil2000cellular} & Multikernel & Cellular Disco: monitor dividido em células e cluster virtual & Balanceamento entre células \\
        \midrule
        \citet{boydwickizer2008corey} & Multikernel & Corey: aplicação controla o compartilhamento das estruturas do kernel & Mecanismos de alocação manuais \\
        \midrule
        \citet{wentzlaff2009fos} & Multikernel & fos: serviços do SO distribuídos por mensagens & Distribuição espacial de serviços \\
        \midrule
        \citet{baumann2009multikernel} & Multikernel & Barrelfish: um kernel por núcleo, mensagens explícitas e estado replicado & Define a arquitetura multikernel \\
        \midrule
        \citet{schupbach2008diversity} & Multikernel & Diversidade de hardware e base de conhecimento & Base de conhecimento para decisões de alocação \\
        \midrule
        \citet{park2012fusedos} & Multi-kernel HPC & FusedOS: kernel leve ao lado de um Linux completo & Partição estática entre kernels \\
        \midrule
        \citet{shimosawa2014ihk} & Multi-kernel HPC & IHK: interface de ciclo de vida & Plano de controle local \\
        \midrule
        \citet{gerofi2016survey} & Multi-kernel HPC & Levantamento e comparação de seis arquiteturas & Particionamento estático versus dinâmico \\
        \midrule
        \citet{beckmann2014pika} & Multikernel & PIKA: pilha de rede dividida em servidores & Balanceamento sem memória compartilhada \\
        \midrule
        \citet{kuo2019multik} & Especialização de kernel & MultiK: visão do código do kernel especializada por aplicação & Especialização por carga \\
        \midrule

% --- Orquestração e balanceamento de carga ---
        \citet{casavant1988taxonomy} & Balanceamento de carga & Taxonomia de escalonamento estático e dinâmico em sistemas distribuídos & Base de classificação \\
        \midrule
        \citet{shivaratri1992load} & Balanceamento de carga & Políticas de transferência, seleção, localização e informação & Estrutura de um balanceador \\
        \midrule
        \citet{mitzenmacher2001power} & Balanceamento de carga & Poder de duas escolhas no balanceamento aleatorizado & Decisões descentralizadas com pouca informação \\
        \midrule
        \citet{blumofe1999work} & Balanceamento de carga & Roubo de tarefas aleatorizado com garantias de desempenho & Balanceamento iniciado pelo receptor \\
        \midrule
        \citet{swetha2025resource} & Orquestração & Revisão de orquestradores de contêineres e lacuna da orquestração por recursos & Lacuna análoga \\
        \midrule
        \citet{lu2025optimized} & Balanceamento de carga & Implantação de containers em dois estágios & Granularidade dividida em dois níveis \\
        \midrule

        \citet{kocoloski2015composition} & Orquestração multikernel & Hobbes: composição de aplicações entre enclaves de um nó & Gerência de enclaves restrita a um nó \\
        \midrule
        \citet{wang2025multikernel} & Orquestração multikernel & Proposta de múltiplas instâncias do Linux via \textit{kexec} & Plataforma-alvo provável do orquestrador \\
        \midrule
        \citet{wang2025lpc} & Orquestração multikernel & Realocação dinâmica de recursos entre instâncias via \textit{device tree} & Mecanismo de redimensionamento de instâncias \\
        \midrule
        \citet{multikernel2025kerf} & Orquestração multikernel & kerf: ciclo de vida e recursos de instâncias de kernel & Ferramenta mais próxima de orquestração, limitada a um único nó \\
        \midrule

% --- Observabilidade e monitoramento ---
        \citet{mansourisamani1993monitoring} & Observabilidade & Modelo funcional de monitoramento de sistemas distribuídos & Base para a coleta de informação \\
        \midrule
        \citet{faseeha2025observability} & Observabilidade & Revisão de observabilidade em microsserviços com \textit{metrics}, \textit{logs} e \textit{traces} & Fundamentos e terminologia da área \\
        \midrule
        \citet{massie2004ganglia} & Observabilidade & Ganglia: monitoramento hierárquico e federado de clusters & Agregação de métricas por nó e entre clusters \\
        \midrule
        \citet{agelastos2014ldms} & Observabilidade multikernel & LDMS: monitoramento contínuo e leve de sistemas de grande escala & Modelo de coleta por nó e agregação em HPC \\
        \midrule
        \citet{kvalnes2015omni} & Observabilidade multikernel & Omni-Kernel: medição e atribuição de todo consumo de recursos & Contabilização por atividade \\
        \midrule
        \citet{west2009coremonitors} & Observabilidade multikernel & Núcleo de hardware dedicado ao monitoramento de desempenho & Núcleo autossuficiente para exportar informações \\
        \midrule

% --- Disponibilidade e isolamento ---
        \citet{ouyang2015pisces} & Isolamento multikernel & Pisces: enclaves com co-kernels leves e isolamento de desempenho & Isolamento de desempenho entre instâncias \\
        \midrule
        \citet{vanthof2022blackbox} & Isolamento & BlackBox: isolar contêineres de um SO não confiável & Imposição de fronteiras a fim de garantir isolamento \\

    \end{longtable}
\end{small}

\section{Considerações Finais}\label{sec:consideracoes-finais}

A análise de todos os trabalhos revela uma trajetória de algumas décadas com uma característica em comum, a gerência
da arquitetura multikernel é sempre feita no escopo de uma única máquina.

É possível evidenciar, principalmente três limitações, que abrem diversas portas para avanço e pesquisa.
Em primeiro lugar, como dito anteriormente, o escopo de gerência é sempre um único nó, no projeto Hive e no Disco é
mencionado clusterização como extensão futura, mas nenhum dos trabalhos analisados de fato as realiza.
Em segundo lugar, os mecanismos frequentemente existem apenas na parte teórica, como o Wax que nunca foi de fato implementado \citep{chapin1995hive},
ou deixam as decisões diretamente à aplicação, ou usuário, como é feito no Corey e no IHK \citep{boydwickizer2008corey,shimosawa2014ihk}.
Em terceiro lugar, as ferramentas de observabilidade tomam como unidade a máquina, o contêiner ou a requisição
\citep{massie2004ganglia,agelastos2014ldms}, e não a instância de kernel.

Entre os trabalhos analisados, nenhum trata a instância de kernel como unidade escalonável, elástica, observável e
recuperável num conjunto amplo de nós.
Sendo assim, existe espaço para a construção de uma camada de orquestração, tomando ideias
de orquestradores de contêineres comerciais da atualidade e estendendo-as para um cluster de nós em ambientes
com a arquitetura multikernel.
Realizando o monitoramento das instâncias, coletando informações dos núcleos e as agregando entre nós, para ser
possível a tomada de decisão da criação de novos jobs, redimensionamento e redistribuição de cargas pelo cluster.
